\chapter{Texturen erstellen}
\label{ch:items/creating-textures}

Texturen\index{textur@Textur} erstellen und optimieren zu können, ist
bei der Entwicklung für Luanti sehr nützlich. Für Pixel-Art-Texturen
gibt es viele hilfreiche Techniken, und wer sie versteht, verbessert die
Qualität seiner Texturen erheblich.

Ausführliche Anleitungen für gute Pixel-Art würden den Rahmen dieses
Buches sprengen; hier geht es nur um die wichtigsten Grundtechniken. Es
gibt viele gute Online-Tutorials\footnote{\url{http://www.photonstorm.com/art/tutorials-art/16x16-pixel-art-tutorial}}
(englisch), die Pixel-Art viel ausführlicher behandeln.

\section{Techniken}\label{techniken}

\subsection{Mit dem Stift arbeiten}\label{mit-dem-stift-arbeiten}

Das Stiftwerkzeug gibt es in den meisten Editoren. Auf die kleinste
Größe eingestellt, bearbeitest du damit ein Pixel nach dem anderen, ohne
andere Teile des Bildes zu verändern. So entstehen klare, scharfe
Texturen ohne ungewollte Unschärfe. Außerdem hast du damit viel
Präzision und Kontrolle.

\subsection{Kacheln}\label{kacheln}

Texturen für Nodes sollten in der Regel kachelbar sein. Das heißt: Setzt
du mehrere Nodes mit derselben Textur nebeneinander, passen die Kanten
nahtlos aneinander.

Passen die Kanten nicht zusammen, sieht das Ergebnis deutlich weniger
ansprechend aus.

\subsection{Transparenz}\label{transparenz}

Transparenz ist wichtig für die Texturen fast aller
Craftitems\index{craftitem@Craftitem} und mancher Nodes, etwa Glas.
Nicht jeder Editor unterstützt Transparenz; achte also darauf, einen
Editor zu wählen, der zu deinen Texturen passt.

\subsection{Farbpaletten}\label{farbpaletten}

Eine einheitliche Farbpalette ist ein einfacher Weg, deine Grafiken
deutlich besser aussehen zu lassen. Verwende am besten eine Palette mit
wenigen Farben, höchstens etwa 32. Fertige Paletten findest du auf
lospec.com\footnote{\url{https://lospec.com/palette-list}}.

\section{Editoren}\label{editoren}

\subsection{MS Paint}\label{ms-paint}

MS Paint ist ein einfacher Editor, der für grundlegendes Texturdesign
ausreichen kann; Transparenz unterstützt er allerdings nicht. Für die
Seiten von Nodes spielt das meist keine Rolle, brauchst du aber
Transparenz, solltest du einen anderen Editor wählen.

\subsection{Aseprite / LibreSprite}\label{aseprite-libresprite}

Aseprite\footnote{\url{https://www.aseprite.org/}} ist ein proprietärer
Pixel-Art-Editor. Er bringt von Haus aus viele nützliche Funktionen mit,
etwa Farbpaletten und Animationswerkzeuge.

LibreSprite\footnote{\url{https://libresprite.github.io/}} ist ein
quelloffener Fork von Aseprite aus der Zeit, bevor Aseprite proprietär
wurde.

\subsection{GIMP}\label{gimp}

GIMP ist in der Luanti-Community weit verbreitet. Die Einarbeitung
dauert recht lange, weil viele Funktionen nicht sofort ersichtlich sind.

Verwende in GIMP das Stiftwerkzeug mit dem Pinsel „Pixel`` und der Größe
1. Außerdem empfiehlt es sich, beim Radierer die Option „Harte Kanten``
zu aktivieren.
